\section{Discussion}\label{sec:discussion}

\paragraph{A small mean difference, but a sign shift} For the late-retracted primary set, the post hoc 90\% interval
rules out absolute mean CD$_5$ differences above 0.00093, about 28\% of the review comparison's effect. Retracted papers
reach the top 5\% of their year about as often as controls, but positive scores are more common (odds ratio 1.23).
Thus, the two primary estimates do not give identical conclusions: the index registers a sign shift without a large
mean difference or unadjusted top-tier excess. An evaluation rewarding high CD$_5$ does not systematically favour
these late-retracted papers, but it does not exclude them either. A separate prospective study found that links to
papers retracted by a fixed cutoff predicted later retraction, albeit weakly as a standalone screen
\citep{arjomandi2026prospective}; \citet{lou2026beyond} reported the retrospective citation-network association.

\paragraph{What the $n_k$ term hides} The papers that cite retracted papers less often also cite their references.
The difference is small, three percentage points in the share of such citing papers or about 0.12 standard deviations
of DI$_{nok}$, but its point estimate is positive in every reason class, though near zero for honest error.
Its interval excludes zero for
misconduct and for other reasons, and the shift in its sign survives adjustment for reference and citation counts.
CD$_5$ registers it only through its sign, because the papers that cite only the references outnumber the citing papers
by more than two orders of magnitude for the median paper. The pattern is in line with the first expectation set out in
the introduction, that fabricated or falsified findings are cited without the literature they built on, but two other
readings fit the data and cannot be separated here: retracted papers may attract citations from outside their immediate
literature, or their reference lists may matter less to the papers that cite them. The data give no support to the
second expectation, that paper-mill products are consolidating. Whatever the reason, variants that remove $n_k$
\citep{wuyan2019solo,bornmann2020convergent} or that adjust for reference length would score retracted papers somewhat
higher than comparable papers. \citet{scancar2026distortion} found citation patterns among suspected paper-mill
articles in highly ranked molecular-cancer journals consistent with citation inflation. This motivates considering
the composition of citing papers, but does not establish why they omit a focal paper's references. The present
analysis does not exclude citing papers that are themselves retracted, so it cannot separate that possibility from
the other explanations.

\paragraph{Reference length and the top tier} At equal reference length, retracted papers reach the top 5\% of
CD$_5$ more often (odds ratio 1.42). The unadjusted comparison, which is what an evaluator sees, shows no
over-representation. The new matched mean-log reference difference, $+0.084$ (0.058--0.109), supports a reference-length
imbalance in the direction that suppresses high scores, consistent with the change after adjustment. The
adjusted estimate shrinks when citations in the window are also held fixed (1.19), but citations are part of the score
itself, so that model may over-adjust. These exploratory models show that the top-tier comparison is sensitive to
reference length; they do not establish a causal effect or validate a preferred adjustment.

\paragraph{Scores that include the retraction} Scores computed over windows that include the retraction differ for two
reasons. Notice records are counted among the citing papers, and more than a third of the papers retracted inside the
window have one. Beyond notices, the papers citing a retracted paper increasingly ignore its references while its
citations fall. This divergence starts before the notice, most clearly for misconduct, which fits critical attention
that precedes a formal retraction, and it steepens after the notice; the comparison is descriptive, not causal. For
users of the index the implication is simple: scores of retracted papers that include the notice year or later years
describe the retraction as much as the paper. They should be computed on a window that ends before the notice, with
notice records excluded, or be flagged.

\paragraph{Honest error} Papers retracted for honest error score slightly lower than their controls on CD$_5$, show no
clear difference on DI$_{nok}$, and their change after the notice shows no sign of an earlier trend. These classes are small (294
and 615 papers in the two windows), so the pattern should be read as a hypothesis: papers later found to contain errors
may build closely on their references, and their errors may come to attention only with the notice.

\paragraph{Relation to the debate on declining disruption} Requiring at least one reference removes almost all
CD$_5 = 1$ records (0.04\% remain), which are the records at the centre of the artefact debate
\citep{holst2024artefacts,holst2026partially,park2026reply}. This study makes no claim about long-run trends; the
yearly means of the reference sample are reported for orientation only. The size of $n_k$ found here, at least 90\% of
the denominator for almost all papers, is a reminder that CD$_5$ values near zero are the rule on a large citation
graph, so that differences between groups appear in thousandths.

\paragraph{Limitations} OpenCitations does not cover all references, and the choice of data source affects disruption
scores \citep{xu2025alternative,crossdi2025}; scores computed here are scores on open data. 83.3\% of Retraction Watch
papers were matched to the graph. Reason classes are a rule-based grouping of Retraction Watch reasons, and many
notices give limited information (the ``other'' class). Because the graph is complete only to about September 2023, the
main analysis ends with papers from 2017 and covers only 276 paper-mill papers; the post hoc three-year window raises
this to 708, but most paper-mill papers were published after 2019. Retractions are a selected and incomplete label of
problematic work: some unreliable papers are never retracted, and the controls may contain them. Several analyses
were added after the main results were known. Their plans were recorded before their runs, but the analyses remain
exploratory. The 0.00093 bound is an interval-based description, not a validated negligible-effect threshold.
Bootstrap resampling clusters by journal--year, not by author; serial retractors can contribute papers across
strata, so uncertainty may be understated if their outcomes are correlated. The diagnostic inputs contain no
author identifiers, preventing an author-cluster or largest-retractor exclusion check.
The primary set is also selected for late retraction: first notices occur in the sixth or a later calendar year after
publication. Its conclusions need not generalize to papers retracted during the scoring window, which are analysed
separately. Notice
records are identified only through PubMed. Because of the earlier trend, the change around the notice cannot be
attributed to the notice alone. Only two members of the index family are analysed.
